\subsubsection{有界算子的谱}\label{subsec:bounded-operator-spectra}

现在研究有界算子的谱的性质. 关于谱的许多最重要的结果 (包括下面关于紧算子的结果所用的结果), 都是利用预解算子的幂级数展开导出的. 现在证明一个基本定理, 它是初等微积分中关于几何级数收敛的结果的类似物.

\begin{Theorem}
\label{thm:neumann-operator-series}
设 $X$ 为 Banach 空间, 并假设 $A\in\mathcal{L}(X)$ 满足 $\|A\|<1$. 那么 $(I-A)^{-1}$ 存在且有界, 而且 $(I-A)^{-1}$ 的以下幂级数展开在算子范数下收敛:
\begin{equation}\label{eq:neumann-operator-series}
(I-A)^{-1}=\sum_{k=0}^\infty A^k=I+A+A^2+\cdots.
\end{equation}

\end{Theorem}

\begingroup
\postdisplaypenalty=10000
\renewcommand{\endProof}{\par\nobreak\hfill$\blacksquare$\par\vspace{0.1cm}}
\begin{Proof}
证明的主要思想是, 若 Banach 空间中的级数\emph{绝对收敛} (即各项范数之和收敛), 则原级数收敛. (这一事实的证明与初等微积分中实数级数的证明相同.) 这里所涉及的 Banach 空间为 $\mathcal{L}(X)$, 且有
\begin{equation}\label{eq:neumann-series-norm-majorant}
\sum_{k=0}^\infty\|A^k\|\leq\sum_{k=0}^\infty\|A\|^k.
\end{equation}
由于 $\|A\|<1$, 右边的几何级数收敛. 因而 \eqref{eq:neumann-operator-series} 右边的级数绝对收敛, 从而收敛. 只需证明其极限确实为 $(I-A)^{-1}$, 而这一步的证明与初等微积分中几何级数的结果基本相同; 即有\[\begin{aligned}
I-A^{k+1}&=(I-A)(I+A+A^2+\cdots+A^k)\\
&=(I+A+A^2+\cdots+A^k)(I-A).
\end{aligned}\]现在, 由于 $\|A\|<1$, 有 $\lim_{k\to\infty}A^{k+1}=0$. 因而
\begin{equation}\label{eq:neumann-series-two-sided-inverse}
I=(I-A)\left(\sum_{k=0}^\infty A^k\right)=\left(\sum_{k=0}^\infty A^k\right)(I-A),
\end{equation}
定理得证.
\end{Proof}
\endgroup

该定理立即给出以下结果, 它表明有界算子 $A$ 的谱 $\sigma(A)$ 位于复平面中的一个有界圆盘内.

\begin{Corollary}
\label{cor:bounded-operator-spectrum-norm-bound}
设 $A\in\mathcal{L}(X)$, 并假设 $\lambda\in\sigma(A)\subset\mathbb{C}$. 那么
\begin{equation}\label{eq:bounded-operator-spectrum-norm-bound}
|\lambda|\leq\|A\|.
\end{equation}

\end{Corollary}

\begingroup
\postdisplaypenalty=10000
\renewcommand{\endProof}{\par\nobreak\hfill$\blacksquare$\par\vspace{0.1cm}}
\begin{Proof}
假设 $|\lambda|>\|A\|$. 那么利用 Theorem~\ref{thm:neumann-operator-series}, 按下式构造预解算子, 即可证明 $\lambda\in\rho(A)$:
\begin{equation}\label{eq:resolvent-series-outside-norm-disk}
R_\lambda(A)=(A-\lambda I)^{-1}=-\dfrac{1}{\lambda}\left(I-\dfrac{1}{\lambda}A\right)^{-1}=-\dfrac{1}{\lambda}\sum_{k=0}^\infty\left(\dfrac{1}{\lambda}A\right)^k.
\end{equation}
这里用到了 $\left\|\dfrac{1}{\lambda}A\right\|<1$ 这一事实. 证明完毕.
\end{Proof}
\endgroup

刚刚已证明有界算子的谱包含于一个圆盘内, 因而自然要问这个圆盘是否最优. 为此给出以下定义.

\begin{Definition}
\label{def:operator-spectral-radius}
从 $X$ 到 $X$ 的算子的\textbf{谱半径}定义为
\begin{equation}\label{eq:operator-spectral-radius-definition}
r_\sigma(A):=\sup_{\lambda\in\sigma(A)}|\lambda|.
\end{equation}

\end{Definition}

因而对于 $A\in\mathcal{L}(X)$, Corollary~\ref{cor:bounded-operator-spectrum-norm-bound} 仅仅表明
\begin{equation}\label{eq:operator-spectral-radius-norm-bound}
r_\sigma(A)\leq\|A\|.
\end{equation}
一般来说, \eqref{eq:operator-spectral-radius-norm-bound} 中的等号不成立, 但对一类称为\emph{正规}算子的算子, 等号成立. 下面的 Problem~\ref{ex:self-adjoint-spectral-radius-equals-norm} 建立自伴算子情形的等式.
